\chapter{价值的重构：在剪切带中栖息}

上一章说明了后果需要被接住，也说明了承受者与决定者可能分离。价值重新分布到哪里，还取决于需求、产权、议价与公共安排。本章区分三件事：某项工作的社会必要性，承担它的人获得的报酬与地位，以及人的尊严和生活意义。三者会互相影响，却没有自动兑现的关系。生成能力充足之后，验证、照护、长期关系与责任可能更加稀缺；这些位置能否获得资源，需要回到制度和分配。历史上的能力外化，提供了考察这种变化的几组材料。

如果智能系统已经能写文案、做总结、翻译、修图、写代码、做方案，普通人的技能还剩下多少重量？如果它还能读文献、提假说、设计实验、跑模拟、分析数据、写论文，科研训练又该如何证明自己的价值？如果它进一步参与数学证明、搜索引理、检查推导，科研中最硬的形式智力劳动是否也开始松动？如果机器人会驾驶、制造、搬运、巡检、手术和照护，那些过去依赖身体、经验和现场判断的工作，还能保留多少人的位置？问题不是某些职业会不会被替代，而是过去用于证明人的训练、经验和主体资格的能力接口，正在被外部系统拆开。

这些问题表面上是在问职业会不会消失，技能会不会贬值，实际上是在问一种更深的东西：过去，人用什么证明自己有价值。会写、会算、会做、会判断、会研究、会证明、会行动，曾经都能支撑一个人在社会中的位置。现在，这些能力正在被外部系统一层层取走、放大、重组。真正让人不安的，不是机器到底还有什么永远做不到，而是当一种能力被外化之后，人该迁移到哪里。

\section{历史的回响：价值漂移的物理学}

这些问题不是第一次被问，答案也不是第一次被改写。第一章已经给出过那条判断：技术每一次外化一种能力，都会切断能力与价值之间的旧绑定，因为被外化的从来不是某一项孤立能力，而是旧社会用来识别训练、经验、资格和责任的接口。这一章要接着往下问的是被切断之后的事：价值落到了哪里，以及它为什么落在那里。三种证明的失效是同时发生的：产出不再天然证明主体，形式不再天然证明经验，能力不再天然证明正当性。把它们拆开可以看到，三者各自依赖一个可以被外部系统接管的环节，而迁移的方向正是由这些环节的位置决定的。

这里说的接口不是抽象说法，它对应一组具体的制度安排。行会用学徒年限确认手艺，科举用考试内容确认资格，大学用文凭确认训练，职业执照用考试和实习期确认准入，学术共同体用发表记录确认可信度。这些安排的共同点是它们把能力的证明落在一个人无法快速获取的过程上：七年学徒期不能压缩成七个月的阅读，住院医师的病例数不能由考试成绩替代，论文的同行评价不能由作者自述替代。接口的作用正是把这种无法快速获取的过程固定成可识别的形式，因此它一旦被外部系统接管，失去的不只是能力的稀缺，还有确认资格的整套程序。\cite{spence1973job}

历史给出的答案从来不是“人不再有价值”，而是价值换了一个挂靠点。这个过程可以沿着几个节点看清：能力被移走的方式不同，旧绑定断裂的方式也不同，而每一次断裂都留下一个相同的问题：被释放出来的那部分价值，接下来要落在哪里。理解这条线不需要完整的技术史，只需要每个节点上最硬的那一步。

\subsection{苏格拉底的恐惧：当智慧失去内化的黏滞}

\subhead{能力外化与旧绑定的断裂}

这里不是要写一部完整技术史，而是抓住几个能力外化的关键节点：记忆、复制、时间刻度、做功、规则执行、符号生成、形式探索、行动能力和生命干预。人类不是第一次经历这种失重。文字让记忆变轻，印刷术让复制变轻，钟表让时间变轻，蒸汽机让肌肉变轻，计算机让规则执行变轻，大模型让符号生成变轻，推理模型和科学代理体让形式探索变轻，具身智能让行动能力开始外化，生命智能让生命干预和生命定义进入技术系统。每一次技术都像一次深层搬迁：它把原本沉积在身体、技艺和身份中的能力移到外部系统里，使旧的价值接口突然松动。但人的价值没有因此消失，而是被迫迁移到更深、更重、更难外包的位置。这九个节点排在一起，是一条价值落点被反复重排的线；重排的结果由谁获得，不取决于技术本身。

节点的顺序不是随意的，它对应抵抗被形式化的能力级别。记忆最先被移走，因为一段话可以被写下来；复制随后被移走，因为一次抄写可以被压成版式；做功被移走得晚一些，因为肌肉必须真的克服重力；规则执行、符号生成、形式探索依次被移走，因为它们分别能被写成流程、写成统计规律、写成可搜索的空间。越是能被写成外部形式的能力，越早离开身体。这个顺序解释了为什么每一次外化都会先打击那些把能力当作身份的人：他们的位置建立在“我能做到别人做不到的事”上，而外部系统恰好把这件事变成了一件便宜的事。

断开的旧绑定有固定的内容。它包含训练的可验证性（这个人确实经过了长时间的准备）、经验的可追溯性（他见过这类对象，并且记得后果）、资格的排他性（不是所有人都能进入这个位置）、责任的归属（出了问题由他承担）。四件事在旧制度里是捆在一起的，一个能力同时证明四件事。外化切断的是捆绳：训练可以外包给模型，经验可以由检索替代，资格可以由订阅获得，责任则在系统之间分散。断口出现的位置，就是价值必须重新落户的位置。

\subhead{记忆的外化与内化的黏滞}

第一个节点是记忆，它外化的时间最早，也最彻底。文字发明以前，记忆不是普通能力，而是秩序本身。法律、神话、谱系、祭祀、交易、禁忌和经验，都必须依赖人的身体保存。能记住、能复述、能传承的人，不只是信息容器，也是共同体时间的承载者。遗忘不是个人失误，而是秩序断裂。文字改变了这一点。它让经验第一次稳定地脱离身体，沉积在符号中。人会死，声音会消失，亲历者会离场，但话语、契约、律令、诗歌、账目和知识可以留下。柏拉图在《斐德若》中借苏格拉底之口警惕文字：它像记忆的药，也可能让人误把外部保存当成内部理解。\cite{plato_phaedrus} 这不是简单反技术，而是看见了第一次\term{价值迁移}：记忆变轻以后，理解变重；文本可以固定答案，却不能替代追问；一个人可以拥有一段话，却未必拥有这段话背后的判断。

苏格拉底的恐惧有一个准确的落点：知识从大脑的生物神经迁移到羊皮卷之后，人对知识的体感不再随文字一起传递。背诵在这里不是低级的存储，而是把内容反复穿过身体、与经验咬合在一起的内化过程，它的黏滞很高，也很贵。文字把这份黏滞降到接近于零，于是“能记住”不再稀缺，稀缺的是“能推演”。文字没有摧毁智慧，它只是换掉了智慧的支点：价值从存储漂移到逻辑构建，从知识的容器漂移到知识的处理器。这场迁移在公元前四世纪就已经完成一半，它留下的问题一直没变：外部保存可以固定答案，却无法替人完成理解。

口传社会的研究给这场迁移提供了一个可以直接测量的对照。Milman Parry 和 Albert Lord 在二十世纪三十年代记录南斯拉夫的口传史诗歌手，发现优秀的歌手并不是把上万行诗句背下来，而是掌握一套格律套语，每次演唱都在现场重新组装；同一位歌手在不同场合唱同一部史诗，长度可以相差几千行。\cite{parry1960singer} 记忆在这种条件下不是仓库，而是一种持续施工的能力，它同时承担着创作、传承和校验。文字把这三件事拆开了：创作归作者，传承归抄本，校验归后来的读者。拆开之后，记忆作为整体能力的黏滞消失了，而理解、校订与解释的黏滞上升了，这也是苏格拉底看到的那一层变化。

这里需要留一个反向的限定，否则这个判断会被读成“文字使人变浅”。书写不只保存了已有的思想，它还制造了新的思想对象。位值制记数和代数符号让运算可以在纸面上一步步展开，证明可以脱离讲述者的记忆被检查，几何作图可以被后人重新走一遍。文本把推理从人脑中搬出来，使它可以被他人停顿、质疑和改写，这件事本身就是推演能力的一次扩张，而不只是记忆能力的一次损失。内化的黏滞被降低，换来的是可公共检查的形式，代价与收益同时成立。

记忆外包之后会发生什么，现代的对照实验给出了一个可以测量的答案。Sparrow、Liu 和 Wegner 在 2011 年发表的一系列实验显示，当被试预期某条信息可以在以后被检索到，他们对这条信息本身的记忆成绩下降，而对“这条信息存在哪里”的记忆成绩上升，研究者把这个现象称为谷歌效应。\cite{sparrow2011google} 这个结果需要按它实际测量的东西来读：它测到的是记忆内容的分布发生了变化，不是记忆能力的净损失，实验也不能说明外部保存一定有害。考试制度的调整提供了一句旁证：允许检索的开卷考试不再考察记忆，转而考察组织、判断与表达，而闭卷考试至今仍在大量场景中使用，理由并不是记忆本身值得尊敬，而是某些领域里的判断必须以足够的内部知识为前提。两份证据指向同一个结构：外化改变的是稀缺的位置，不是人的总量。

记忆变轻只是第一次，复制变轻的幅度更大，而它动的不是个人的理解力，是权威本身。当文本可以被大量复制，谁有权解释文本就取代了谁拥有文本，成为新的稀缺位置。这是一次和记忆完全不同的迁移：前一次迁移发生在个人的能力结构里，这一次迁移发生在共同体的权力结构里。

\subsection{古腾堡的降维：解释权垄断的黏滞坍缩}

\subhead{复制的降维与解释权的坍缩}

印刷术进一步降低了复制的重量。手抄时代，文本本身稀缺，拥有书、抄写书、保存书，就是权力。印刷术让文本大量扩散，使知识从少数空间中涌出，进入更广阔的社会循环。Eisenstein 把印刷术视为近代欧洲知识秩序转变的重要力量，因为它改变的不只是复制速度，也包括校勘、传播、争论和权威形成方式。\cite{eisenstein1979printing} 可是复制变轻以后，人的价值没有消失，而是从复制迁移到筛选、解释、校订、组织信任和建立公共判断。文本越多，判断越稀缺；复制越便宜，信任越昂贵。印刷没有让人无关，反而让“谁解释、谁校验、谁建立可信版本”变得更重要。

这一步降维的对象不是文字，而是复制本身。手抄受制于抄写员的物理速度，一位抄写员一天能完成的字数有限，经文因此稀缺，解释权因此可以被垄断在少数人手里：谁掌握抄本，谁就掌握通往真理的入口。当复制的边际成本趋近于零，信息以接近超流体的方式扩散，稀缺的位置从经文转移到诠释，1517 年马丁·路德把《九十五条论纲》送上印刷机，几年之内传遍欧洲，靠的不是教会的准许，而是复制已经不再昂贵。诠释随即通胀，旧权威在低黏滞环境里失去支点，恐慌的本质不是信息过载，而是解释权换了主人。价值也从“拥有信息”漂移到“定义信仰”，再漂移到为信息建立可信版本，这条漂移线在往后的每一次外化里都会重演。

复制的降维还有一个不显眼的后果：它让校勘和版本比对第一次成为可以规模化的技术工作。同一部书的不同版本可以被并排比较，异文可以被系统整理，错误可以在重印中修正，作者可以在生前反复修订自己的著作。这些工作在过去依赖少数抄写室的传统，后来变成出版机构、图书馆和学术团体的常规职能。复制变便宜的同时，确认哪一个版本可信的成本上升了，于是新的位置出现在校订者、编纂者和出版者一边。稀缺没有消失，它只是从“有没有文本”换到了“文本是否可信”。

复制和生成变便宜，既可能扩展知识，也会增加筛选与建立可信版本的需求。确认本身也可以被工具加速。需要观察的是整条链上哪种功能相对稀缺，以及承担这项功能的组织能否得到资源。

\subhead{规则执行与符号生成}

复制变轻以后，被移走的能力不再停留在复制一侧，而是转向了可被写成流程的那部分心智劳动。规则的重量无法用吨和瓦特计量，迁移的结构却相同。计算机又把这种迁移推进到心智劳动。它降低的不是全部思考，而是可形式化、可规则化、可程序化过程的成本。手算、账务、表格、数据存储、检索、排序、重复执行和确定性流程，被机器大规模吸收。一个人算得再快，也不可能以这种方式重新赢回价值。于是人的位置迁移到更上游：定义问题，建立模型，选择变量，理解语境，判断输出是否对应现实需求。计算机让执行变轻，却让问题定义变重。Simon 讨论人工物和有限理性时已经指出，人的智能并不是脱离环境的纯粹求解，而是在目标、约束、信息和组织结构中设计可行路径。\cite{simon1969sciences} 当执行被外包，真正稀缺的不是执行本身，而是知道什么值得执行、为什么执行、哪些问题可以形式化、模型边界在哪里、结果是否应该进入现实。

会计行业在几十年里完成过一次完整的演示。手工记账、试算平衡、报表汇总这些环节被电子表格和财务软件吸收之后，记账员的数量下降了，但审计、税务筹划、内部控制和财务解释的位置没有减少，反而因为数据量变大而增加。变化不在于工作量，而在于能力证明的方向：能把账做平不再证明什么，能判断一笔交易应该如何被处理、能否经得起事后检查，才构成资格。这条路径和计算机在其他确定性流程密集行业里造成的路径一致，只是它发生得更早，也更容易观察。

大模型的冲击更深，因为它进入了人最习惯用来证明自己的区域：语言、表达、解释、代码、证明、方案、风格、总结、计划和初级判断。蒸汽机替代肌肉，计算机替代规则执行，大模型则开始替代大量符号产出。它不是一个更聪明的写作工具，而是文字之后最重要的符号事件之一。文字让经验脱离身体，印刷让经验大规模复制，计算机让符号按规则执行，大模型让已经外化的人类经验重新组合、回应、生成和代理化。外部记忆不再只是保存和等待检索，它开始主动说话、主动解释、主动计划、主动调用工具。

这个变化已经超出“模型能答题、写代码、过考试”的阶段。GPT-5.5 的系统卡把它定位到复杂现实工作，包括写代码、在线研究、分析信息、创建文档和表格，并在多工具环境中推进任务。\cite{openai2026gpt55systemcard} o3 和 o4-mini 的系统卡则把推理模型与网页浏览、Python、图像和文件分析、自动化、文件搜索和记忆等工具能力结合起来。\cite{openai2025o3o4minisystemcard} Deep Research 和 ChatGPT Agent 进一步说明，研究、检索、网页交互、代码运行、数据分析和文件生成正在合流。\cite{openai2025deepresearchsystemcard,openai2025chatgptagentsystemcard} 这些系统共同指向一个趋势：大模型正在从符号生成器转向工具化、代理化的任务系统。人的价值问题也因此不能停留在“谁写得更好”。如果模型能写、能查、能算、能跑代码、能调用工具、能生成报告，那么人的价值必须迁移到更深的位置：谁确认信息，谁定义问题，谁判断结果可以进入现实，谁接管失败，谁承担后果。

但镜子不是世界。大模型可以生成知识的外形，却不自动拥有经验；可以生成证明的外形，却未必真正穿过规则；可以生成关怀的外形，却没有承受关系；可以生成行动方案，却没有接触后果。过去，一段判断通常暗示背后有某个说话者、某段经历、某种责任；现在，形式可以在没有经历、没有验证、没有承受的情况下大量出现。形式和经验之间的旧联系被削弱了。大模型让文明的外部记忆会说话，也让“说出来”本身失去原有重量。

\subhead{生成之后的验证序列}

外部记忆开始说话之后，价值的位置需要重新安排。验证可以按介质差异分成形式、语义、因果、工程、生命长期和社会责任六层，那是按对象分类；换一个角度，按后果的深度排列，同一条链可以读成一个逐级上升的序列，每一级都比上一级更慢、更贵、更不容易被替换。候选与凭据之间的距离，正是价值要重新沉淀的地方。

所以，大模型之后，人的第一层价值不是继续生产更多符号，而是让符号重新接回现实。生成越便宜，验证越稀缺；候选越多，判断越重要。答案不是知识，解释不是理解，方案不是行动，模拟不是后果。智能时代不是知识自然变多，而是候选爆炸。只有穿过验证金字塔的候选，才可能成为可依赖的结果。这里所谓价值，不只是劳动市场中的交换价值，也包括知识共同体中的可信性、制度系统中的授权性，以及关系和生命世界中的承受性。AI 可以降低产出的成本，却不会自动生成这些价值的正当性。

这座金字塔的层级并不复杂：候选生成 $\to$ 表征确认 $\to$ 推演确认 $\to$ 行动确认 $\to$ 系统嵌入 $\to$ 长期承受。候选生成回答“能不能说出来”；表征确认回答“对象是否存在、来源是否可靠”；推演确认回答“证明、代码、计算和因果链是否成立”；行动确认回答“能否进入真实环境，错误能否停止、回滚和追责”；系统嵌入回答“组织、平台、医院、道路、工厂和制度能否接住”；长期承受回答“能否经受时间、规模、异常、滥用和代际后果的检验”。智能系统当然也可以参与验证，甚至在许多局部环节上比人更稳定；但这只说明部分验证也会被外化，并不意味着验证的正当性自动完成。智能可以生成判断，也可以验证判断，却不能自动授予判断进入现实的正当性。

链条一旦展开，就能看出各段的迁移速度并不一样。越靠前的环节越容易被自动化替代，因为它的输入和输出都可以被表述清楚；越靠后的环节越慢，因为它的判据要在时间中才能获得，而且判据的持有者必须自己承担判断错误的代价。科研同时包含最硬的形式环节和最重的对象环节，先变的是前一半，这不是因为后一半更复杂，而是因为后一半的判据不在系统内部。这一层区别，把下一步的讨论推向了两端：规则空间里的搜索，和真实对象上的接触。

\subsection{钟表与蒸汽机：被规训的时间与被解放的肌肉}

\subhead{钟表先规训了时间}

在蒸汽机之前，机械钟表已经先动过一次手，它外化的不是肌肉，而是时间。自然节律下的工作随日升日落展开，小时是模糊的，赶工看天色和季节，人对节奏保留着相当大的自主。钟楼和怀表把时间切成均匀、可计量、可分割的刻度，工厂汽笛把刻度变成纪律，准时成为美德，迟到成为过失，效率第一次可以被计时比较。这一层外化不产出任何实物，却重排了所有人的生活，工业秩序的地基因此不是煤和铁，而是被规定的时间；人在这里失去的不是力气，而是对节奏的支配。

蒸汽机改变的不是符号，而是能量。它让人和动物的肌肉不再是主要动力来源。过去，体力、耐力、速度和直接做功能支撑大量社会位置；蒸汽机之后，力量开始沉积在机器、燃料、铁路、工厂、城市和资本中。人的身体不再是做功的核心单元，社会开始围绕机器功率重组。旧手艺失重，旧劳动被重新纪律化，但人没有消失。人的价值从肌肉迁移到控制、组织、维护、调度和责任。机器越有力量，事故越严重；功率越大，标准、维修、培训、调度、保险和责任越重要。工业革命不是简单让人从劳动中解放出来，而是让身体做功失去旧价值，把人的位置推向机器系统的设计、管理、协调和承责。\cite{landes1969unbound}

钟表和蒸汽机合起来完成的是同一次迁移：人从做功者变成控制者，价值锚点从克服物理阻力漂移到管理复杂系统。控制者的位置比做功者的位置更抽象，它要求的是把机器、时刻表和陌生人组织起来的能力，因此也更难被一次技术升级抹掉。这几层外化有一个共同点，被移走的能力都在外部世界里留下了可见的、可计量的载体，因此旧的位置还能被重新找到。机器开始吸收规则、符号和形式探索之后，迁移的方向就不再这么直观。

\subsection{被外化的形式与行动}

\subhead{形式探索的外化}

科研也会因此重新定位，其中数学和理论研究只是最尖锐的一端。这里需要区分三个层级：符号生成是说出形式，形式探索是在规则空间中产生新结构，验证则是让新结构被规则、共同体或现实对象接住。过去，能写出证明、能提出构造、能搜索算法，本身就是高强度训练的证明；现在，这些接口也开始松动。AlphaGeometry 在奥林匹克几何问题中把神经生成与符号演绎结合起来：神经模型提出辅助构造，符号引擎负责推演和检查，在无人工示范的情况下合成大量几何定理与证明。\cite{trinh2024alphageometry} 这个例子说明，可靠证明不是“像证明的文字”，而是形式生成与规则确认的结合。

更强的信号出现在 2026 年：OpenAI 模型给出反例，推翻了与 Erdős 1946 年平面单位距离问题相关的核心猜想方向；随后 Alon、Bloom、Gowers、Litt、Sawin、Shankar、Tsimerman、Wang 和 Wood 给出一份短的、消化过的、人类验证版本，并对结果进行反思。\cite{openai2026unitdistancedisproof,alon2026remarks} 这里必须说得精确：这不是说单位距离问题已经被一句话“解决”，也不是说数学共同体可以被省略；它说明，连最硬的形式智力接口也开始松动。模型能够提出真正改变问题图像的构造，但这个构造要成为数学知识，仍然需要人类数学家检查、压缩、解释、命名和嵌入理论。

AlphaEvolve 展示了另一种形式探索外化方式。它不是简单让模型写一段证明，而是把大模型生成、代码修改、自动评估和演化搜索结合起来，使候选算法在反馈中迭代改进。白皮书称，它在若干数学、计算和基础设施优化问题中产生了有用结果；其中最适合作为例子的，是它在 4$\times$4 复矩阵乘法这一特定设置中找到使用 48 次标量乘法的程序，改进了长期未被突破的结果。\cite{novikov2025alphaevolve} 这个案例不应被泛化成“AI 已经解决算法发现”，它更准确地说明：在可形式化、可评估、可迭代的搜索空间里，候选生成和候选筛选可以被机器闭环加速。人的价值因此不是消失，而是迁移到对问题、评价函数、验证标准和共同体准入的判断上：什么问题值得搜索，评价器是否正确，候选是否只是过拟合评估器，结果是否可解释、可证明、可迁移、可进入知识体系。这类外化之所以还停在形式一侧，是因为评价器本身可以被写成代码；一旦评价必须由对象给出，迁移就会换一种形态。

\subhead{前沿难题中的构造与确认}

上面两个案例合起来指向一个更大的问题：形式探索的外化能推到多硬的问题上。数学中被单独标出的那一组难题（黎曼猜想、Birch 和 Swinnerton-Dyer 猜想、纳维-斯托克斯方程解的存在性与光滑性一类）有特殊地位，不只因为难，而因为它们的判定标准极其干净：命题要么被证明，要么被推翻，中间没有解释空间，也没有行业惯例可以商量。任何外部系统如果在这类问题上产生实质推进，它碰到的是本书全书讨论的那条最硬的线。

按介质响应谱的说法，这里要问的是阻抗匹配：同一束智能输入进入数学，与进入实验科学，被吸收的方式完全不同。数学一侧的阻抗低，因为它的验证器是形式化的，一个候选对象只要合法就能被独立检查；实验一侧的阻抗高，因为验证器是对象本身，而对象不接受任何形式的替代。这个差别解释了为什么形式探索会先在这一侧松动，也解释了为什么它的松动不能外推。

同一个图像在二〇二六年由陶哲轩以另一套词汇独立提出，而它出现的位置比第三章更靠内。第三章的失配发生在智能输出与外部介质之间，源端是模型，负载端是原子、生命或组织；这里的两端都在知识生产的内部，源端是生成，负载端是消化。他把数学工作拆成生成、验证、消化三件事，消化指的是理解一个结果、把它放进上下文、把它讲清楚。历史上三者的难度相当，消化作为副产品自然发生，因此共同体奖励的是生成与验证。生成能力提高之后，三者的速率分开了：生成最快，验证因为形式化也跟着加快了一部分，消化几乎没有变快。他把这个分离直接称为阻抗失配，并说数学正从证明稀缺进入证明过剩，而它的文化和基础设施还没有适应这个转变。\cite{tao2026aiviews}

生成与消化的分离，先表现为两种能力的比值。设$g>0$为单位时间新增候选的数量，$d$为同一口径下可供消化的能力，能力覆盖比记为
\begin{equation}
\begin{split}
\rho&=\min\!\left\{1,\frac{d}{g}\right\},\\
\left.\frac{\mathrm{d}\rho}{\mathrm{d}t}\right|_{0<d/g<1}&<0
\quad\Longleftrightarrow\quad
\frac{\dot g}{g}>\frac{\dot d}{d}
\end{split}
\label{eq:digestion-share}
\end{equation}
$g$与$d$均以候选数每单位时间计，$\rho$无量纲。在$0<d<g$的区间，覆盖比下降当且仅当生成能力的相对增长率高于消化能力；$d\ge g$时，覆盖比按1计。它衡量的是当期能力与新增负荷的关系，实际已消化份额还要考虑积压、筛选顺序与候选差异。生成从数月缩短到一天，而成文只缩短少量时间时，两类能力的增长会明显分离。图\ref{fig:proof-abundance}示意这种负荷变化。

\begin{figure}[ht]
\centering
% >>>>>>>>>> figs/fig-proof-abundance.tex >>>>>>>>>>
% 图 6-4  证明过剩：生成、验证与消化的速率分离
% 左图 (a)：三条吞吐曲线。纵轴是相对量，不标数值刻度——三条线的间距才是要读的东西，
% 标上数值反而会诱导读者去做它支持不了的比较。横轴中点用底色分出两个时期。
% 生成与消化之间的阴影不是「两条线的高低差」，而是被生成出来却没有被消化的那部分产出，
% 位置相当于第三章反射系数里没有进入负载的那份功率。验证线画在中间，是因为形式化
% 同时加快了一部分验证，但没有加快消化。曲线用折线不用平滑曲线：横轴是示意性的时期，
% 平滑化会假装中间存在一个连续可导的过渡。
% 右图 (b)：同一时期文献构成的翻转。两段柱子总长相同，只换上下顺序，
% 强调变化的是比例而不是总量。百分比只取两个数，避免把示意做成统计。
% 图内文字一律短标签，不放整句：句子进图注。
\begin{tikzpicture}[x=1cm, y=1cm, font=\small,
  ax/.style={-{Stealth[length=2.0mm]}, ink, line width=0.6pt}
]
\useasboundingbox (-0.55,3.95) rectangle (12.90,-1.42);

% ==================== (a) 三条吞吐曲线 ====================
% 两个时期的底色
\fill[mist]    (0.60,0.62) rectangle (4.30,3.62);
\fill[rust!8]  (4.30,0.62) rectangle (7.15,3.62);
\node[font=\scriptsize, text=slate]    at (2.45,3.44) {证明稀缺};
\node[font=\scriptsize, text=rust!85]  at (5.72,3.44) {证明过剩};

% 失配阴影：上边界沿生成线正走，下边界沿消化线倒走
\fill[rust!20]
  (4.30,1.27) (5.10,1.60) (5.90,2.20) (6.70,3.05) (7.15,3.52)
  (7.15,1.31) (6.70,1.27) (5.90,1.20) (5.10,1.12) (4.30,1.06)
  -- cycle;

% 三条曲线
\draw[deepblue, line width=1.9pt]
  (0.60,0.98) -- (2.45,1.10) -- (4.30,1.27) -- (5.10,1.60)
  -- (5.90,2.20) -- (6.70,3.05) -- (7.15,3.52);
\draw[forest!90, line width=1.5pt]
  (0.60,0.92) -- (2.45,1.02) -- (4.30,1.14) -- (5.10,1.34)
  -- (5.90,1.62) -- (6.70,1.90) -- (7.15,2.05);
\draw[terra, line width=1.5pt]
  (0.60,0.88) -- (2.45,0.97) -- (4.30,1.06) -- (5.10,1.12)
  -- (5.90,1.20) -- (6.70,1.27) -- (7.15,1.31);

% 图例：放在左上角没有曲线经过的空区
\draw[deepblue,  line width=1.9pt] (0.90,2.86) -- (1.42,2.86);
\node[font=\scriptsize, text=deepblue, anchor=west] at (1.50,2.86) {生成};
\draw[forest!90, line width=1.5pt] (0.90,2.48) -- (1.42,2.48);
\node[font=\scriptsize, text=forest!90, anchor=west] at (1.50,2.48) {验证};
\draw[terra, line width=1.5pt] (0.90,2.10) -- (1.42,2.10);
\node[font=\scriptsize, text=terra, anchor=west] at (1.50,2.10) {消化};

% 失配标注：生成线斜穿整块阴影，文字放进带内任何位置都会压线。
% 改为在带内竖一支双向箭头，标签退到消化线下方的空区，与箭头同列，不需要引线。
\draw[rust!80, line width=0.9pt, {Stealth[length=1.6mm]}-{Stealth[length=1.6mm]}]
      (6.40,1.30) -- (6.40,2.68);
\node[font=\scriptsize, text=rust!90, anchor=north] at (6.40,0.98) {阻抗失配};

% 坐标轴
\draw[ax] (0.20,0.62) -- (7.62,0.62);
\draw[ax] (0.60,0.62) -- (0.60,3.78);
\node[font=\scriptsize, text=ink, anchor=north] at (3.95,-0.02) {时间};
\node[font=\scriptsize, text=ink, rotate=90, anchor=south] at (0.22,2.20)
      {相对吞吐（示意）};

% ==================== (b) 文献构成的翻转 ====================
\node[font=\scriptsize, text=ink, anchor=south] at (10.36,3.34) {消化覆盖（示意）};

% 两根横向堆叠条
\fill[forest!62] (8.62,2.42) rectangle (10.78,2.97);
\fill[rust!45]   (10.78,2.42) rectangle (11.02,2.97);
\fill[rust!45]   (8.62,1.42) rectangle (10.78,1.97);
\fill[forest!62] (10.78,1.42) rectangle (11.02,1.97);

\node[font=\scriptsize, text=ink, anchor=east] at (8.50,2.70) {提升前};
\node[font=\scriptsize, text=ink, anchor=east] at (8.50,1.70) {提升后};

\node[font=\scriptsize, text=white, anchor=center] at (9.70,2.70) {已消化为主};
\node[font=\scriptsize, text=white, anchor=center] at (9.70,1.70) {未消化为主};
\node[font=\scriptsize, text=rust!90,  anchor=west] at (11.12,2.70) {未消化};
\node[font=\scriptsize, text=forest!90, anchor=west] at (11.12,1.70) {已消化};

% 两个速率：数据而不是句子，所以留在图内
\node[font=\scriptsize, text=ink, anchor=north west, align=left] at (8.62,0.92)
      {生成：6 个月 $\rightarrow$ 1 天\\[2pt]
       成文：1 个月 $\rightarrow$ 3 周};

% ==================== 子图编号，一律放下面 ====================
\node[font=\small, text=ink] at (3.90,-0.86) {(a)};
\node[font=\small, text=ink] at (10.36,-0.86) {(b)};
\end{tikzpicture}

% <<<<<<<<<< figs/fig-proof-abundance.tex <<<<<<<<<<

\caption{证明过剩的机制。(a) 生成、验证、消化的相对吞吐随时间的变化，纵轴为示意刻度，不标数值；两条时期底色标出转向的位置。生成线与消化线之间的距离是被生成出来却没有被消化的那部分产出，表示新增候选与消化能力的失配；验证线位于中间，是因为形式化加快了一部分验证，却没有加快消化。(b) 同一时期文献构成的翻转，两段柱子总长相同，只换上下顺序，变化的是比例而不是总量。两个速率数字取自陶哲轩的算例，柱形只示意覆盖程度，不表示统计份额。\cite{tao2026aiviews}}
\label{fig:proof-abundance}
\end{figure}

反例型与搜索型进展，改变知识的方式有所不同。反例要求给出违反命题条件的对象，搜索型工作则在明确评价器下改进候选。AlphaEvolve的矩阵乘法工作属于后者。可检查的对象为自动搜索提供了反馈，但构造和验证之间的成本关系随问题变化。结果成为可用知识，还需要准确陈述、独立检查与对其适用范围的理解。

雅可比猜想与纳维-斯托克斯问题，提供了两种需要检查具体构造的例子。雅可比猜想自1939年提出，讨论满足相应雅可比条件的多项式映射是否可逆；2026年7月，Levent Alp\"oge公布了一个三维反例，并将发现归于Claude Fable模型。\cite{alpoge2026jacobian} 随后的解释工作把多项式构造放回代数几何结构中，帮助理解反例为何成立。\cite{tao2026jacobiandigestion} 发现、检查与解释可以由不同参与者完成，三者的成本也需要分别考察。

纳维-斯托克斯存在性与光滑性属于千禧年问题。2026年9月8日，OpenAI公布了带光滑外力的三维流动在有限时间形成奇性的证明稿与Lean形式化，并将其作为相应问题版本的解答。\cite{openai2026navierstokes} 这项工作展示了候选构造、分析推导和机器检查的组合。其数学地位需要按命题条件、证明内容、独立核查与共同体审议逐项判断。

黎曼猜想中的零点比例结果，属于向完整命题推进的部分定理。\cite{anthropic2026riemann,lamzouri2026} Birch和Swinnerton-Dyer猜想则涉及椭圆曲线算术与其$L$函数之间的一般关系。部分证明、一般证明与反例构造要求不同的对象和论证。构造可能缩短核查，也可能本身极难寻找；任何一类进展先出现，都需要由具体问题与方法解释。

结果的核查与发现过程的追溯也需要分开。一个反例可以接受独立计算，产生它的搜索过程仍可能没有完整公开；分析证明与形式化证书则各自依赖所采用的命题、规则与实现。纳维-斯托克斯发布前后出现的贡献与数据使用争议，又把问题扩展到来源、保密与署名。相关报道记录了不同参与者的说法，\cite{axios2026mathcredit,techcrunch2026buckmaster,venturebeat2026codexdata} 判断这些争议需要相应记录，不能由结果是否通过检查单独解决。

还有一层值得记录：这一系列事件是从一条后来被否定的传言开始的。此前的说法是 Anthropic 已经解决了两个千禧年难题，OpenAI 在其后启动自己的尝试，先让智能体评估全部六个未解问题，再集中到纳维-斯托克斯上。传闻被证伪，但由它触发的工作留下了结果。这意味着在这类问题上，公开信息本身的信噪比已经成为一个变量。一条不实的传言也能在几天内调动起数千亿 token 的算力。被改变的因此不只是检查的优先级，还有检查的对象选得对不对。

把这些放在一起，可以得出一句可检验的话：在这类问题上，外部系统真正的贡献不是给出答案，而是把一个方向从“被认为是对的”变成“必须重新检查”。这个动作在科学史上一直存在，过去由人完成，现在有一部分可以由系统完成，而且更快。它改变的不是知识的存量，而是检查的优先级。而前面那个分离给出了另一条限制：检查的优先级可以提高得很快，把检查结果变成可用的知识却不能同样快，两者的差额就是这一节留下的待结余额。

\subhead{旧评价体系的失效}

形式探索一旦被部分外化，最先承压的不是知识本身，而是学术界用来筛选知识的那套装置。这套装置的功能用本书的词汇说很清楚：它是一个凭据生成器。同行评议、期刊分级、引用计数、职称与终身教职这一整套安排，生产的是“这篇结论已经被有资格的人检查过”这一凭据，而凭据之所以可信，是因为它由不可压缩的过程生成：评审要花时间，复现要花时间，一个结论被共同体接受要等下一代人重新用到它。

生成能力的上升同时冲击这条链的两端。在产出端，投稿量、候选证明、候选实验方案的供给以数量级上升，而评议能力由人力时间决定，几乎不动；在筛选端，判据本来依赖一个隐含前提（一个人不可能在短时间内产出这么多高质量工作），而这个前提正在失效。两端同时受力，结果是整套装置的通过率与信噪比一起下降，而它自身的调整速度受制于人的工作节奏。这不是学者个人的道德问题，是结构问题：一个靠产出稀缺来证明资格的制度，在产出不再稀缺时会失去分辨力。

式 \eqref{eq:digestion-share} 里的 $\rho$ 在旧安排下之所以不成问题，是因为生成与消化被绑在同一件作品上：一个人解出一个问题，他就必须把它写清楚，否则这篇东西发表不出去。绑定一旦松开，$\rho$ 的下降就直接落到评议端。陶哲轩把这一段的后果写成一条漂移链：优先权的标尺从期刊出版日期漂到预印本的时间戳，再从预印本漂到社交媒体上“谁先让 AI 把提示词跑出来”的竞赛，而这个竞赛奖励的恰好是省掉验证、文献回顾和成文。他提出的替代方案是让标尺回到解释：优先权属于第一个能公开把结果讲清楚的人，形式包括讲座，也包括一篇说明性的文章或一次录了像的报告，并且这个解释性版本是预印本与形式化证书的补充而不是替代。二〇二六年创刊的期刊 Mathematical Discourse 以研究报告的录像作为发表形式，可以看作这条路线的一个早期落点。\cite{tao2026aiviews,tao2026discourse}

同一个思路还推到了一条更硬的规范上。他把证明的完整生命周期拆成五段：生成、验证、成文、发表、经典化；经典化指的是一个结果变成教科书里的常识，能被后续工作反复调用。按这个标准，用 AI 生成并验证一个证明、随后对写成文章、讲成报告、送进评议都没有兴趣，等于把一个结果丢在中间阶段就离开，而接手的人要付出的是别人的成本。他提出的规范是：任何以一个证明的生成主张优先权的人，应当尽力把它推进到至少发表阶段；在他自己也承认这个承诺并非总能兑现之后，他又提出应当有一个正式的对应机制，让走不完流程的人可以把一个证明明确“挂出领养”，交给另一组作者继续，而不是把这类情况留给完全临时的处理。\cite{tao2026aiviews}

验证一侧也出现了同类分工。陶哲轩参与推动的 Palomar 是一个 Lean 形式化数学的登记库，它要对付的是生成过剩带来的一个具体麻烦：一个 Lean 代码库通过了类型检查，并不等于它证明的就是它声称的那个命题。它设了三道检查：代码能否通过类型检查，有没有通过添加公理之类的取巧，形式化的陈述与它声称的非形式陈述在语义上是否一致。第一道是机械的，由 Lean 的比对工具完成；第二道和第三道由大语言模型执行。Palomar 自己写明，这三道检查加起来仍然远达不到人类同行评议在原创性、重要性与准确性上给出的判断，并明确声明它不是一份经过同行评议的期刊。\cite{tao2026palomar}

登记库将验证拆成类型检查、假设审查和陈述对应检查，使不同环节可以分别评估。模型承担中间检查后，还要记录其依据、复核覆盖与遗漏类型。形式证书、语义审查、同行评议和责任安排共同构成确认结构。哪一层成为瓶颈，要由实际吞吐与错误记录判断。

对比两边的说法可以看清这件事的性质。产业一侧的表述是能力与门槛。二〇二六年九月三日，OpenAI 在新模型发布的同一天宣告“欢迎来到通用人工智能时代”，其总裁称它最终可能被视为通用人工智能本身，英伟达的首席执行官在发布之后也直接说通用人工智能已经到来；公司公布的是基准分数、智能体规模、训练算力这些量，并称这次训练首次在得州的一个站点动用了十万块以上的 GPU。\cite{openai2026pathastra,axios2026agi,huang2026agi} 专业共同体一侧的表述是制度与后果：陶哲轩在同年九月的一份声明里写道，数学正在迅速接近那条他不希望走上的路径；在他判断中危害的主要来源已经变了：二〇二六年以前是“用了却不验证”，现在是以基准为驱动、不加区分地榨取；他还说，产业内与他看法相同的人大多已经离开或被边缘化，并把纳维-斯托克斯这一事件称为那种“不计社会价值、只求做出极强自主系统”的竞赛最直观的实例。\cite{tao2026statement}

这两种说法并不互相矛盾，它们量的不是同一件事。前者量的是能力上限，后者量的是能力落地时谁承担代价。本书前面几章反复给出的结论也在这里：能力上限可以单独发布，代价不能单独发布，它必须落到某个位置上；当两边的叙述不一致时，通常不是有一方说谎，而是一方在报分子，另一方在报分母。

评价制度的调整需要比较不同凭据。独立复现能暴露部分错误，代价是额外的试验和人力；被实际使用可以提供长期反馈，也可能使错误先进入后果链；署名、记录与责任安排则有助于追溯决定。三者各自回答不同问题，制度需要组合这些证据，并防止入口被少数机构垄断。

上面几段的材料给第三条路补了一个注。承诺与领养机制、以及把优先权移向“第一个讲清楚的人”，本质上都是把凭据重新挂到责任上：主张优先权的人要为结果的后续处境负责，而不是只为它的出生负责。陶哲轩把这条路线放在一个更大的判断之下。二〇二六年他写道，数学面临的不是一个数学论证的危机，而是一个价值与实践的危机，共同体只有几个月时间重新说明“做一个数学家意味着什么”，如果不自己回答这些问题，答案会由一家技术公司给出，或者由财务激励决定。\cite{tao2026icm,tao2026nature} 这个判断把前面几节的结论收拢到一句话上：当生成一侧的产能不再稀缺，制度必须找一处不可加速的地方来挂靠，而需要持续投入的一项工作，是有人为结果的后续使用、修订和错误处理负责。

三条路都不完整，它们之间也不互相排他。可以确定的只有一件事：新安排不会是旧安排加上更强的工具，因为旧安排的分辨力恰恰来自产出稀缺。这与前面关于凭据的判断一致：当一类凭据的生成过程可以被加速时，它的证明力就下降；证明力下降之后，制度必须重新找一处不可加速的地方来挂靠。至于那处地方在哪里，本书只能给出判据，给不出答案。

\subhead{实验与代码中的外化}

科研的另一半工作在对象一侧，它落成两件事：让候选接触样品和仪器，以及让代码进入真实系统。实验型科研又比形式探索多一层重量。它不只包括读论文、写代码和提出假说，还包括设计实验、制备样品、操控仪器、观察现象、处理误差、重复验证和承受失败。The AI Scientist-v2 声称可以自主提出假说、设计并执行实验、分析和可视化数据，并撰写论文，其中一篇完全 AI 生成的论文达到 ICLR workshop 接受阈值。\cite{yamada2025aiscientistv2} 这个案例必须降调使用：它显示科研流程的若干形式环节已经可以被代理化，但不能证明科学共同体的判断、实验对象接触和长期验证可以被取消。它真正揭示的是，研究者的价值不会继续停留在“能不能完成论文流水线”，而会迁移到问题重要性、对象接触、变量选择、测量可信度、失败解释、伦理边界和长期验证责任。

Virtual Lab 则提供了更贴近生命科学的强案例。它把多个 LLM agent 组织成“虚拟实验室”，由 LLM principal investigator 协调不同 scientist agent，并在人类研究者高层反馈下设计 SARS-CoV-2 变体纳米抗体。\cite{swanson2025virtuallab} 这个案例的价值不是“AI 独立做生物学”，而是说明实验设计、候选筛选和跨学科讨论可以被代理化；但最后仍要进入蛋白模型、实验对象、测量体系和人类判断。科学不是纯 bit 的闭环。候选可以在模型空间里大量生成，科学结果却必须穿过对象、仪器、误差、复现和共同体。

代码位于符号和现实之间。代码看起来像文本，但它能运行、调用接口、写入系统、改变业务流程。GPT-4.1 在 SWE-bench Verified 上被报告完成 54.6\% 的任务，GPT-5-Codex 进一步面向 agentic coding，在真实编码任务中通过强化学习训练，并迭代运行测试直到通过。\cite{openai2025gpt41,openai2025gpt5codex} 这些结果说明，软件生产中的代码生成、补丁候选、测试修复和仓库导航正在被外化。但软件工程的价值没有因此只剩下“写代码”。真正沉重的是需求定义、架构边界、依赖治理、安全审查、测试覆盖、部署责任和长期维护。代码生成越快，代码确认越重要；补丁越容易生成，系统债务越容易堆积。人的价值从“能不能写出一段代码”迁移到“能不能让代码在系统中长期可靠地存在”。到这里，确认工作仍然发生在系统内部；外化如果再向外推一层，被取代的对象就从符号换成动作。

\subhead{行动与生命的外化}

动作和生命过程的后果不可回滚，这一层的外化和前面几层不同。实验型科研中的自动化操作，主要还是为了验证候选假设；具身智能会把问题推到更外层。智能不再只是为了实验而操作对象，而是开始在道路、工厂、医院、家庭和公共空间中直接行动。RT-2 已经把视觉-语言模型与机器人轨迹数据共同训练，使互联网尺度的语义知识以动作 token 的形式进入机器人控制。\cite{zitkovich2023rt2} Gemini Robotics 和 GR00T N1 则把 VLA 模型、人形机器人基础模型和真实机器人轨迹、人类视频、合成数据连接起来，显示数字领域的通用能力正在尝试向物理智能体迁移。\cite{geminirobotics2025,nvidia2025gr00t} 但 RoboCasa365 这样的基准又从反面提醒我们：通用行动不是会说话就能落地。365 个日常厨房任务、2500 个环境、600 多小时人类演示和 1600 多小时合成演示，恰恰说明行动能力必须被场景、任务、数据和评估体系压实。\cite{nasiriany2026robocasa365}

机器一旦进入现实，就要面对接触、磨损、延迟、遮挡、滑移、误抓、失稳、损坏、停机、他人身体和场景变化。生命体的身体经过长期演化，许多复杂性被封装在形态、反射和习惯中；机器没有自然继承这笔摊销，所以行动能力越强，越需要验证、维护、停机、追溯和承责结构。因此，在具身智能时代，人的价值不只是验证知识，还包括裁量行动。自动驾驶能不能上路，不取决于它在演示视频里多流畅，而取决于极端场景、事故责任、接管机制、道路规则、保险结构和公众可接受风险是否被安排；医疗机器人能不能进入手术室，不取决于它是否能完成某个动作，而取决于失败时谁接管、谁记录、谁解释、谁承担损失。行动不是输出的延长，而是进入后果链的准入事件。机器可以参与行动，但后果不能悬空。把后果固定到一个可以操作的位置上，需要登记、日志、审计、停机权和保险，这些东西不直接产出效率，却构成行动能够发生的前提。

如果说大模型让符号产出失去重量，推理模型和科学代理体让形式探索开始外化，具身智能让行动能力开始外化，那么生命智能会把问题推到更深处：智能不再只是替人说、替人算、替人做，而是开始参与定义身体、疾病、正常、风险和生命路径。AlphaFold 3 已经把结构预测从单个蛋白质推进到蛋白质、核酸、小分子、离子和修饰残基等复合物的相互作用预测。\cite{abramson2024alphafold3} AI 药物发现、蛋白设计、基因风险评分和多模态医疗模型，也正在把生命过程中的候选生成、筛选、分类和风险预测推进到更高密度。\cite{fritzsche2023ethical}

但生命智能不是直接把生命变轻，而是把生命中的分类、筛选、风险和优化目标提前到模型空间。分子结构预测让候选更容易出现，风险评分让未来疾病概率更容易被分类，多模态医疗模型让症状、影像、基因、病历和生活方式更容易被合并为判断。问题随之改变：什么算健康？什么算正常？什么差异应被保留？什么痛苦可以消除？什么风险不能转嫁？谁有权为尚未出生的人选择未来？谁有权定义“更好的人”？当生命进入设计空间，价值问题就不再是人有没有用，而是什么样的人、痛苦、差异和风险有资格存在。

生命智能最危险的地方，不在于它生成了多少候选分子，而在于它可能把某些分类提前写进人的未来：谁被提前视为高风险，谁被提前视为低价值，谁被提前要求治疗、筛选、优化或排除。智能可以提供定义的候选，却不能独自完成定义的授权。分类一旦进入制度，就会改变谁被视为病人，谁被视为风险，谁被视为缺陷，谁被要求治疗、筛选或优化。即使未来智能在流行病学、基因组学、脑科学和风险预测上远超人类专家，问题也不会消失，只会后移：谁授权这种标准进入制度，谁允许它改变人的分类、待遇、保险、教育、生殖和医疗机会，谁承受它长期错误的后果？如果授权程序本身也被智能系统深度渗透，就更需要把患者经验、公共讨论、法律程序和长期后果检验保留下来，否则所谓“更好的人”，就可能变成技术系统自我强化的分类结果。走到这一步，外化已经不只是能力问题，而是授权问题；价值迁移也就不再只是个人处境，而是制度处境。

\subsection{能力外化与价值迁移}

\subhead{迁移的社会条件}

记忆、复制、计时和规则执行在历史上多次被外化，技术降低了相应操作的成本，也改变了训练、职业与解释权的价值。新的困难可能转向校准、组织和分配，也可能被后续技术继续降低。人的位置随这些关系重排。判断这次迁移，需要分别考察哪些工作仍被需要、哪些人获得收益，以及哪些代价没有进入核算。

剩下的问题是分配：新的位置由谁获得，又由谁决定。从文字到智能时代，人类历史呈现出一条深线：人不断把自身的一部分能力外化出去，每一次外化都让某种旧能力变轻，也让人退到更深的位置上。人不是越来越没有价值，而是不能再用旧位置证明自己。

但高黏滞区并不天然属于人。验证可以被平台化，准入可以被资本化，责任可以被合同化，定义权也可以被专家系统和制度机器重新分配。人的价值迁移不是因为人天然占据更高位置，而是因为符号、形式、行动和生命一旦进入现实，就必然需要某种可追问、可申诉、可拒绝、可承担的接口。如果这些接口被少数平台、机构或算法系统垄断，人并不会自动升格，只会从生产者变成被分类、被评分、被调参的对象。所谓人的价值迁移，不是一条温和的上升曲线，而是一场围绕验证权、准入权、承责权和定义权的社会博弈。

这也意味着，大多数人的价值迁移不可能只靠个人努力完成。文字出现后，不是每个记忆者都自然成为解释者；印刷术出现后，不是每个抄写者都自然进入大学和出版制度；蒸汽机出现后，也不是每个手工业者都自然变成工程师和管理者。智能时代同样如此。不是每个被替代的文员都会自然成为验证专家，不是每个被挤压的劳动者都会进入决策层。人的新位置也可能存在于地方经验、现场判断、照护关系、使用反馈、拒绝权、申诉权和公共监督中。关键不只是培养少数高端裁量者，而是让更多人在新的技术结构中仍有表达真实经验、拒绝错误目标、参与后果分配的通道。否则，价值迁移就会变成少数人的上升通道和多数人的功能性失重。

迁徙的失败不只是失去位置，还包括时间上的错位：新的位置不会在旧的位置消失的同一时刻出现，中间的间隔由被替代者自己承担。劳动力市场的研究把这条机制的轮廓描得比较清楚：自动化同时产生替代效应与互补效应，前者压缩可被标准化的任务，后者提高与机器配合的工作的生产率，两者在同一时期发生，但落点不是同一群人。\cite{autor2015why} 美国制造业就业人数在 1979 年前后达到约 1960 万人的峰值，到 2010 年前后降到约 1160 万人，而这一段时间里全国的就业总量仍在增长。\cite{bls2024employment} 总量与分布的分离，是这类迁徙最容易被宏观叙事抹掉的部分；中间期在个人尺度上并不短，一项技能的持有者通常在四十岁以后才遇到替代压力，而重新形成一项有溢价的能力需要数年，这段时间里的损失由个人承担，收益要到下一代人身上才显现。

失败并不只有失去位置这一种形态。价值迁移的另一种失败，不是人彻底无用，而是人被安置在一个过于顺滑的环境中：系统替人生成、筛选、推荐、执行和安慰，人的判断、拒绝、等待和承受逐渐失去锻炼空间。第五章讨论过的停滞与局部最优，说的就是这个结构在系统一侧的样子：拒绝失去对象，系统停在一个舒适的局部最优上。没有失败，判断无法校准；没有等待，欲望无法变厚；没有拒绝，边界无法形成；没有外部反作用，系统只会越来越贴合已有偏好。人需要的不是痛苦本身，而是可以失败、可以否定、可以反悔、可以重新定义目标的空间。没有这种空间，主体不会成熟，只会被环境持续调参。

人提供的不是单纯的扰动、反馈或偏好信号，而是带有历史、意图、关系和后果承受的差异性反应。系统当然可以把人的选择转化为数据，把人的反馈转化为参数，把人的行为转化为优化信号；但这种转化不能穷尽人的意义。只有会受伤、会记忆、会拒绝、会反悔、会承担后果的主体，才能把系统从自洽回路中拉回开放世界。人不是智能系统的低级输入，也不是数据燃料；人是意义和后果仍然发生的地方。
\begin{figure}[ht]
\centering
% >>>>>>>>>> figs/fig-value-drift.tex >>>>>>>>>>
% 图 6-2  价值漂移的阻力迁徙
% 横轴按「可重复、可计量、可拆成步骤」的程度排列，从左到右依次变难形式化。
% 两条折线是两个时点的阻力地形：虚线是尚未被技术接手时的样子，
% 实线是今天的形状。两条线之间的面积就是被填平的阻力，
% 它集中在左段；右段两条线几乎重合，说明那一侧的阻力没有被技术动过。
% 因此「价值守恒」在这里说的不是分量不变，而是填平一块洼地会在高处抬出新位置。
% 折线用直线段而非光滑曲线：横轴是十个离散类别，圆滑化会掩盖这个事实。
% 坐标一律写成字面数字，不经 \xdef 累加——累加版本在本文件里静默失效过一次。
% 十个类别的横坐标依次为 0.65 起、间隔 1.30。
\begin{tikzpicture}[x=1cm, y=1cm, font=\small,
  ax/.style={-{Stealth[length=2.0mm]}, ink, line width=0.6pt}
]
\useasboundingbox (-0.35,1.45) rectangle (12.95,-5.60);

% ---------- 被填平的阻力（先铺底）：上边界正走，下边界倒走 ----------
\fill[rust!14]
  (0.65,-0.55) (1.95,-0.43) (3.25,-0.18) (4.55,-0.80) (5.85,-1.05)
  (7.15,-1.43) (8.45,-1.30) (9.75,-0.93) (11.05,-0.80) (12.35,-0.55)
  (12.35,-0.11) (11.05,-0.61) (9.75,-0.99) (8.45,-1.36) (7.15,-1.49)
  (5.85,-2.86) (4.55,-2.80) (3.25,-2.93) (1.95,-3.05) (0.65,-2.99)
  -- cycle;

% ---------- 技术接手之前：虚线 ----------
\draw[slate!75, line width=1.4pt, dashed]
  (0.65,-0.55) -- (1.95,-0.43) -- (3.25,-0.18) -- (4.55,-0.80) -- (5.85,-1.05)
  -- (7.15,-1.43) -- (8.45,-1.30) -- (9.75,-0.93) -- (11.05,-0.80) -- (12.35,-0.55);
% ---------- 今天：实线 ----------
\draw[forest!90, line width=1.9pt]
  (0.65,-2.99) -- (1.95,-3.05) -- (3.25,-2.93) -- (4.55,-2.80) -- (5.85,-2.86)
  -- (7.15,-1.49) -- (8.45,-1.36) -- (9.75,-0.99) -- (11.05,-0.61) -- (12.35,-0.11);

% ---------- 图例：放在左下角的大片空区 ----------
\node[font=\scriptsize, text=slate!85, anchor=west] at (0.85,-1.30)
      {虚线：技术接手之前};
\node[font=\scriptsize, text=forest!90, anchor=west] at (0.85,-1.75)
      {实线：今天};

% ---------- 被填平的区域 ----------
\node[font=\scriptsize, text=rust!85, anchor=center] at (4.85,-1.76)
      {被填平的阻力};

% ---------- 压缩前沿 ----------
\draw[ink!55, line width=0.6pt, densely dotted] (6.50,-3.30) -- (6.50,-0.05);
\node[font=\scriptsize, text=ink, anchor=west] at (6.62,-0.28) {压缩前沿};

% ---------- 坐标轴与十个类别 ----------
\draw[ax] (0.20,-3.30) -- (12.95,-3.30);
\foreach \xx/\nm in {0.65/记忆,1.95/复制,3.25/体力,4.55/时间刻度,5.85/规则执行,
                     7.15/验证,8.45/准入,9.75/承责,11.05/定义,12.35/承受} {
  \node[font=\scriptsize, text=ink, anchor=north] at (\xx,-3.42) {\nm};
}
\node[font=\scriptsize, text=slate, anchor=north] at (3.25,-4.02)
      {可重复、可计量、可拆成步骤};
\node[font=\scriptsize, text=slate, anchor=north] at (9.75,-4.02)
      {不可形式化、后果由承受者承担};

% ---------- 两个流向 ----------
\draw[rust!80, line width=1.0pt, -{Stealth[length=1.8mm]}] (6.11,0.85) -- (0.85,0.85);
\node[font=\scriptsize, text=rust!85, anchor=south] at (3.64,0.93)
      {技术优先流向低黏滞区};
\draw[forest!90, line width=1.0pt, -{Stealth[length=1.8mm]}] (6.89,0.85) -- (12.65,0.85);
\node[font=\scriptsize, text=forest!90, anchor=south] at (9.88,0.93)
      {价值最终沉淀在高黏滞区};

% ---------- 底部说明 ----------
\node[font=\scriptsize, text=ink, anchor=north west, align=left]
      at (0.00,-4.62)
      {人的位置不是被填平，而是被推到更难压缩的地带：右段两条线几乎重合，\\[2pt]
       说明那里的阻力没有被技术动过，它也就成了价值重新沉淀的地方。};
\end{tikzpicture}

% <<<<<<<<<< figs/fig-value-drift.tex <<<<<<<<<<

\caption{价值漂移的阻力迁徙。横轴按形式化的难易排列，从左到右依次变难；两条折线是两个时点的阻力地形。要读的是两线之间的间距沿横轴如何分布：它给出的是阻力迁徙的方向，而不是迁徙的总量。}
\label{fig:value-drift}
\end{figure}

\subhead{不可完全封闭的开口}

所以，人的价值不能建立在“机器永远不能成为主体”这种断言上。未来智能会发展成什么样，不能轻易封死。更深的判断是：任何智能系统，只要进入符号、形式、行动和生命，就不只是在提高效率，而是在重写价值接口。符号需要现实校准，形式需要验证共同体，行动需要责任结构，生命优化需要授权程序和承受主体。\key{技术优先流向低黏滞区，价值最终沉淀在高黏滞区。}生成、复制、搜索和初步执行会越来越轻；验证、准入、承责、定义、长期关系、身体接触、公共讨论和不可回滚的承诺会越来越重。

人的价值不是不可替代的产出，而是不可完全封闭的现实开口。智能越能替代人的能力，人越不能把价值押在能力本身。人的价值必须迁移到更深的边界：让符号接回现实，让形式接回验证，让行动接回责任，让生命接回不可外包的承受。这四次要接回的东西都不能被生成，因为它们各自带着一种无法被复制的代价：验证要押上判断者自己的可信性，准入要押上可以停机的能力，承责要押上可以赔付的资源，定义要押上可以被追问的授权。价值迁移的方向因此不是一串更强的能力，而是一组更贵的承诺；代价能否构成可信性，决定了这些位置能不能真的接住价值。

这条漂移线给出一组位置，却没有说明这些位置由什么构成。要让判断成立、承诺有效、责任可追溯，系统必须先有边界，而边界不是靠声明获得的。接下来要处理的就是这件事：边界在物理上如何被定义，以及一个东西要凭什么才能证明自己确实付出了代价。
